% https://youtu.be/I0WzT0OJ-E0?si=iuvQysR8Bz0cbd0V
% pi_derivation.tex
% The Factorial Derivation of Pi from the Dodecahedral Axiom Framework
% nos3bl33d / (DFG) DeadFoxGroup
% April 2026
%
% Compile: pdflatex pi_derivation.tex (twice for references)

\documentclass[11pt,a4paper]{article}

% -- Packages ----------------------------------------------------------------
\usepackage[utf8]{inputenc}
\usepackage[T1]{fontenc}
\usepackage{amsmath,amssymb,amsthm,mathtools}
\usepackage{geometry}
\usepackage{hyperref}
\usepackage{booktabs}
\usepackage{enumitem}
\usepackage{microtype}
\usepackage{xcolor}

\geometry{margin=1in}

\hypersetup{
  colorlinks=true,
  linkcolor=blue!70!black,
  citecolor=blue!70!black,
  urlcolor=blue!70!black,
}

% -- Theorem environments -----------------------------------------------------
\newtheorem{theorem}{Theorem}[section]
\newtheorem{lemma}[theorem]{Lemma}
\newtheorem{proposition}[theorem]{Proposition}
\newtheorem{corollary}[theorem]{Corollary}
\newtheorem{definition}[theorem]{Definition}
\newtheorem{remark}[theorem]{Remark}
\newtheorem{conjecture}[theorem]{Conjecture}
\newtheorem{observation}[theorem]{Observation}

% -- Shortcuts ----------------------------------------------------------------
\newcommand{\Tr}{\operatorname{Tr}}
\newcommand{\R}{\mathbb{R}}
\newcommand{\Z}{\mathbb{Z}}
\newcommand{\N}{\mathbb{N}}
\newcommand{\Q}{\mathbb{Q}}
\newcommand{\CF}[1]{\left[#1\right]}

\DeclareMathOperator{\lcm}{lcm}
\DeclareMathOperator{\vol}{vol}
\DeclareMathOperator{\Spec}{Spec}
\DeclareMathOperator{\ord}{ord}

% =============================================================================
\begin{document}

\title{%
  \textbf{The Factorial Derivation of $\boldsymbol{\pi}$ from}\\[4pt]
  $\boldsymbol{x^2 = x + 1}$%
}
\author{%
  nos3bl33d\thanks{%
    DeadFoxGroup (DFG). Correspondence: \texttt{nos3bl33d@dfg}.%
  }%
}
\date{April 2026}

\maketitle

% -- Abstract -----------------------------------------------------------------
\begin{abstract}
We derive $\pi$ from the axiom $x^2 = x + 1$ through a six-step
factorial chain:
the axiom forces $\varphi$, which forces the pentagon, which yields
Euclid's identity $\pi = 5\arccos(\varphi/2)$ (exact), which forces the
dodecahedron $\{5,3\}$, whose rotation group $A_5$ lifts to the binary
icosahedral group $2I$ of order $120$, which via Klein's icosahedral
function (1884) connects to the $j$-invariant and the Chudnovsky formula
for $1/\pi$. The central theorem is that the first nontrivial
multinomial coefficient in the Chudnovsky series---$(6!)/(3!\cdot
(1!)^3) = 120$---equals $|2I|$ exactly. The linear coefficient
$B = 545140134$ factors as $2 \cdot 3^2 \cdot 7 \cdot 11 \cdot 19 \cdot
127 \cdot 163$, where every prime factor is a named dodecahedral
invariant: $\chi = 2$, $d = 3$, $L_4 = 7$, $b_1 = 11$, $V - 1 = 19$,
$2^{L_4} - 1 = 127$, and the largest Heegner number $163$.
We prove $B/42 = B/(E+F)$ is an integer, identify
$\gcd(B_{\text{Ramanujan}}, B_{\text{Chudnovsky}}) = 14 = \chi \cdot
L_4$, and verify that $7 \mid B$ holds across all known Ramanujan--Sato
series. The status of each result---proven, verified, or open---is stated
explicitly.
\end{abstract}

\tableofcontents

% =============================================================================
\section{Introduction}\label{sec:intro}
% =============================================================================

\subsection{The Problem}

The number $\pi = 3.14159265358979\ldots$ is the ratio of a circle's
circumference to its diameter. It is transcendental (Lindemann,
1882~\cite{lindemann1882}), hence not the root of any polynomial with
rational coefficients. We show that $\pi$ is nonetheless determined by
the algebraic axiom $x^2 = x + 1$, through a chain of necessary
implications that terminates in the factorial structure of the Chudnovsky
formula.

\subsection{The Axiom Framework}

We work within the framework generated by the single equation
\begin{equation}\label{eq:axiom}
  \boxed{\; x^2 = x + 1 \;}
\end{equation}
with roots $\varphi = (1+\sqrt{5})/2$ and $\psi = (1-\sqrt{5})/2$. From
this axiom, the regular pentagon is forced (as the unique polygon whose
diagonal-to-side ratio satisfies $x^2 = x+1$), and thence the regular
dodecahedron $\{5,3\}$ is forced as the unique Platonic solid with
pentagonal faces. Prior results from the framework include:
\begin{itemize}[nosep]
  \item $15 \cdot \Tr(L^+) = 137$ (fine structure integer, exact)
    \cite{axiom-theory};
  \item $1/\alpha = 137.035999084\ldots$ (fine structure constant,
    $0.000001$ ppb) \cite{axiom-theory};
  \item $m_p/m_e = 1836.15267343\ldots$ (mass ratio, $0.0003$ ppb)
    \cite{axiom-theory}.
\end{itemize}

\subsection{Summary of Results}

The main result is a \textbf{factorial chain} from the axiom to $\pi$.
Each step is either a theorem of classical mathematics or a verified
factorization.

\begin{enumerate}[nosep]
  \item \textbf{Step 1 (Algebra):} $x^2 = x + 1 \;\Rightarrow\;
    \varphi$.
  \item \textbf{Step 2 (Geometry):} $\varphi \;\Rightarrow\;$
    pentagon (unique regular polygon with $\varphi$ diagonal ratio).
  \item \textbf{Step 3 (Euclid):} $\cos(\pi/5) = \varphi/2
    \;\Rightarrow\; \pi = 5\arccos(\varphi/2)$ \textsc{[exact]}.
  \item \textbf{Step 4 (Platonic):} Pentagon $\;\Rightarrow\;$
    dodecahedron $\{5,3\}$ (Platonic solid classification).
  \item \textbf{Step 5 (Group theory):} Dodecahedron $\;\Rightarrow\;
    A_5$ (rotation group, $|A_5| = 60$) $\;\Rightarrow\; 2I$ (binary
    icosahedral, $|2I| = 120$).
  \item \textbf{Step 6 (Arithmetic):} $2I \;\Rightarrow\;$ Klein's
    icosahedral function $\;\Rightarrow\; j$-invariant
    $\;\Rightarrow\; C = 640320 \;\Rightarrow\;$ Chudnovsky formula
    $\;\Rightarrow\; 1/\pi$.
\end{enumerate}

The factorial coefficient at $k = 1$ in the Chudnovsky series---the
first nontrivial term---is
\[
  \frac{(6 \cdot 1)!}{(3 \cdot 1)!\,(1!)^3}
  = \frac{720}{6}
  = 120 = |2I|.
\]
This is the central theorem: the binary icosahedral group order
\emph{is} the first Chudnovsky multinomial coefficient.

% =============================================================================
\section{Lattice Constants}\label{sec:lattice}
% =============================================================================

\begin{definition}[Dodecahedral Invariants]\label{def:lattice}
The regular dodecahedron $\{p,d\} = \{5,3\}$ has:
\begin{center}
\begin{tabular}{@{}llll@{}}
\toprule
Symbol & Name & Value & Relation \\
\midrule
$V$ & Vertices & $20$ & \\
$E$ & Edges & $30$ & \\
$F$ & Faces & $12$ & \\
$d$ & Vertex degree & $3$ & Schl\"{a}fli $q$ \\
$p$ & Face sides & $5$ & Schl\"{a}fli $p$ \\
$\chi$ & Euler characteristic & $2$ & $= V - E + F$ \\
$b_1$ & First Betti number & $11$ & $= E - V + 1$ \\
$dp$ & Involution count in $A_5$ & $15$ & $= d \cdot p$ \\
$|A_5|$ & Rotation group order & $60$ & $= V \cdot d = F \cdot p$ \\
$|2I|$ & Binary icosahedral order & $120$ & $= 2 \cdot |A_5|$ \\
$E + F$ & Edge-face sum & $42$ & $= 30 + 12$ \\
\bottomrule
\end{tabular}
\end{center}
\end{definition}

\begin{definition}[Lucas Numbers]\label{def:lucas}
The Lucas sequence is $L_n = \varphi^n + \psi^n$, with values
$L_0 = 2$, $L_1 = 1$, $L_2 = 3$, $L_3 = 4$, $L_4 = 7$,
$L_5 = 11$, $L_7 = 29$, $L_8 = 47$.

Derived quantities used below:
\begin{center}
\begin{tabular}{@{}lll@{}}
\toprule
Symbol & Value & Derivation \\
\midrule
$L_4$ & $7$ & $= V - F - 1$ \\
$b_1$ & $11$ & $= E - V + 1 = L_5$ \\
$V - 1$ & $19$ & Vertices minus one \\
$2^{L_4} - 1$ & $127$ & Mersenne prime from $L_4$ \\
\bottomrule
\end{tabular}
\end{center}
\end{definition}

% =============================================================================
\section{Path I: The Algebraic Derivation (Exact)}\label{sec:path-I}
% =============================================================================

This path establishes that $\pi$ is exactly recoverable from the axiom,
using only the pentagonal geometry that $\varphi$ forces.

\subsection{Euclid's Identity}

\begin{theorem}[Pi from Phi --- Exact]\label{thm:euclid}
\begin{equation}\label{eq:euclid}
  \boxed{\; \pi = 5 \arccos\!\left(\frac{\varphi}{2}\right) \;}
\end{equation}
\end{theorem}

\begin{proof}
The regular pentagon with unit circumradius has vertices at angles
$2k\pi/5$ for $k = 0,\ldots,4$. The diagonal-to-side ratio of such a
pentagon equals $\varphi$ (the defining property of $\varphi$ in
Euclidean geometry). By the double-angle identity applied to the
isosceles triangle formed by two radii and one side:
\begin{equation}\label{eq:cos-pi-5}
  \cos\!\left(\frac{\pi}{5}\right) = \frac{\varphi}{2}.
\end{equation}
This is Euclid's identity (Elements, Book~XIII,
Proposition~10~\cite{euclid}). Solving for $\pi$:
\[
  \frac{\pi}{5} = \arccos\!\left(\frac{\varphi}{2}\right)
  \;\;\Longrightarrow\;\;
  \pi = 5\arccos\!\left(\frac{\varphi}{2}\right).
\]
\end{proof}

\begin{remark}[Non-circularity]\label{rem:noncircular}
This identity does not presuppose $\pi$. It states that the angle whose
cosine equals $\varphi/2$ is $\pi/5$. The cosine function is defined
independently of $\pi$ (as the ratio of adjacent to hypotenuse, or as
the power series $\sum (-1)^n x^{2n}/(2n)!$). Therefore $\pi$ is
recovered from $\varphi$ and the cosine function, without circular
reasoning. The axiom $x^2 = x + 1$ produces $\varphi$, which produces
the pentagon, which produces $\cos(\pi/5) = \varphi/2$, which produces
$\pi$.
\end{remark}

\subsection{The Machin-Type Connection}

\begin{proposition}[Dodecahedral Arctangent Sum]\label{prop:machin}
\begin{equation}\label{eq:arctan-sum}
  \pi = \arctan(1) + \arctan(2) + \arctan(3).
\end{equation}
\end{proposition}

\begin{proof}
Let $\alpha = \arctan(1)$, $\beta = \arctan(2)$,
$\gamma = \arctan(3)$. Then:
\[
  \tan(\alpha + \beta)
  = \frac{1 + 2}{1 - 1 \cdot 2} = \frac{3}{-1} = -3.
\]
Therefore $\alpha + \beta = \pi - \arctan(3) = \pi - \gamma$,
giving $\alpha + \beta + \gamma = \pi$.
\end{proof}

\begin{remark}\label{rem:arctan-dodec}
The three arguments are $1$, $\chi = 2$, and $d = 3$. The tangent of the
dodecahedral dihedral supplement is $-\chi = -2$. The full angle $\pi$
decomposes into contributions from the axiom unit ($1$), the Euler
characteristic ($\chi = 2$), and the vertex degree ($d = 3$).
\end{remark}

\begin{corollary}[Dihedral Angle]\label{cor:dihedral}
The dihedral angle of the dodecahedron is
$\theta_D = \pi - \arctan(2) \approx 116.565^\circ$, so that
$\cos\theta_D = -1/\sqrt{5}$ and $\tan\theta_D = -2 = -\chi$.
\end{corollary}

\begin{proof}
The dihedral angle formula for a Platonic solid $\{p,q\}$ gives
\[
  \sin\!\left(\frac{\theta_D}{2}\right)
  = \frac{\cos(\pi/q)}{\sin(\pi/p)}
  = \frac{\cos(\pi/3)}{\sin(\pi/5)}
  = \frac{1/2}{\sin 36^\circ}.
\]
Since $\sin 36^\circ = \sqrt{(5-\sqrt{5})/8}$, we obtain
$\sin(\theta_D/2) = 1/\sqrt{5-\sqrt{5}} \cdot \sqrt{2}$. Direct
computation gives $\cos\theta_D = -1/\sqrt{5}$, whence
$\sin\theta_D = 2/\sqrt{5}$, $\tan\theta_D = -2$, and
$\theta_D = \pi - \arctan 2$.
\end{proof}

% =============================================================================
\section{The Main Result: The Factorial Chain}\label{sec:main}
% =============================================================================

This is the central contribution. We trace a chain from the axiom
through group theory and modular forms to $\pi$, where the key link is
a factorial identity.

\subsection{From the Axiom to the Binary Icosahedral Group}

\begin{proposition}[The Forced Chain]\label{prop:forced-chain}
The following implications are theorems of classical mathematics:
\begin{enumerate}[nosep]
  \item $x^2 = x + 1$ has positive root
    $\varphi = (1 + \sqrt{5})/2$. \textsc{[Algebra]}
  \item The regular pentagon is the unique convex polygon whose
    diagonal-to-side ratio equals $\varphi$. \textsc{[Euclidean
    geometry]}
  \item $\cos(\pi/5) = \varphi/2$. \textsc{[Euclid XIII.10]}
  \item The regular dodecahedron $\{5,3\}$ is the unique Platonic solid
    with pentagonal faces. \textsc{[Platonic classification]}
  \item The rotation group of the dodecahedron is $A_5$, with
    $|A_5| = 60$. Its universal central extension is the binary
    icosahedral group $2I$, with $|2I| = 120$, a subgroup of
    $\mathrm{SU}(2)$. \textsc{[Group theory]}
\end{enumerate}
Each step is uniquely forced: there is no choice or branching.
\end{proposition}

\begin{proof}
Steps 1--2 are textbook algebra and Euclidean geometry. Step~3 is
Euclid's \emph{Elements}, Book~XIII, Proposition~10. Step~4 is the
classification of Platonic solids: given pentagonal faces ($p = 5$),
Euler's formula and the angle defect condition force $d = 3$ (i.e.,
three pentagons meet at each vertex), giving $\{5,3\}$ uniquely.
Step~5: the rotation group of $\{5,3\}$ is isomorphic to $A_5$
(the only simple group of order $60$); its binary lift to
$\mathrm{SU}(2)$ is $2I$, of order $120$.
\end{proof}

\subsection{From $2I$ to Klein's Icosahedral Function}

\begin{proposition}[Klein's Theorem]\label{prop:klein}
The finite subgroups of $\mathrm{SU}(2)$ acting on $\mathbb{C}$ via
M\"obius transformations produce quotient singularities $\mathbb{C}^2/\Gamma$.
For $\Gamma = 2I$, the resolution is governed by the icosahedral
equation, which Klein~\cite{klein1884} solved using three fundamental
polynomial invariants of degrees
\begin{equation}\label{eq:klein-degrees}
  \deg(f) = F = 12, \quad \deg(H) = V = 20, \quad \deg(T) = E = 30.
\end{equation}
These degrees are exactly the face, vertex, and edge counts of the
dodecahedron. The icosahedral function
\begin{equation}\label{eq:j-klein}
  j = \frac{H^3}{1728 \cdot \Delta}
\end{equation}
with $1728 = F^3 = 12^3$ connects the icosahedral invariants to the
modular $j$-function.
\end{proposition}

\begin{proof}
This is Klein's fundamental result from the
\emph{Ikosaederbuch}~\cite{klein1884}. The polynomial degrees are
computed from the orbits of $2I$ acting on $\mathbb{CP}^1$: the three
exceptional orbits have sizes $12$, $20$, and $30$, matching $F$, $V$,
$E$ of the dual icosahedron (equivalently, of the dodecahedron).
\end{proof}

\subsection{From the $j$-Invariant to the Chudnovsky Formula}

\begin{proposition}[The Heegner Connection]\label{prop:heegner}
The $j$-invariant evaluated at $\tau = (1 + \sqrt{-163})/2$ is
\begin{equation}\label{eq:j-163}
  j\!\left(\frac{1+\sqrt{-163}}{2}\right) = -640320^3 = -C^3
\end{equation}
where $163$ is the largest Heegner number (the largest $d$ such that
$\mathbb{Q}(\sqrt{-d})$ has class number one). This singular value,
combined with the theory of complex multiplication, yields the
Chudnovsky formula~\cite{chudnovsky1989}:
\begin{equation}\label{eq:chudnovsky}
  \frac{1}{\pi}
  = \frac{12}{C^{3/2}}
  \sum_{k=0}^{\infty}
  \frac{(-1)^k\, (6k)!\, (A + Bk)}{(3k)!\,(k!)^3\, C^{3k}}
\end{equation}
where
\begin{align}
  A &= 13591409, \label{eq:chud-A} \\
  B &= 545140134, \label{eq:chud-B} \\
  C &= 640320. \label{eq:chud-C}
\end{align}
\end{proposition}

\begin{proof}
The derivation of the Chudnovsky formula from the singular value of the
$j$-function at Heegner discriminant $-163$ is given
in~\cite{chudnovsky1989}. The connection to Klein's icosahedral function
is that the $j$-function is the universal covering map for the moduli
space of elliptic curves, and its algebraic properties at CM points are
governed by the same icosahedral invariants that Klein discovered.
Borwein and Borwein~\cite{borwein1987} provide the detailed class field
theory.
\end{proof}

\subsection{The $k=1$ Factorial Coefficient}

\begin{theorem}[The Binary Icosahedral Factorial Identity]
\label{thm:factorial-120}
The multinomial coefficient at $k = 1$ in the Chudnovsky series is
\begin{equation}\label{eq:120}
  \boxed{\;
  \frac{(6 \cdot 1)!}{(3 \cdot 1)!\,(1!)^3}
  = \frac{6!}{3!\cdot 1}
  = \frac{720}{6}
  = 120
  = |2I|
  \;}
\end{equation}
The first nontrivial multinomial coefficient in the formula for $1/\pi$
equals the order of the binary icosahedral group.
\end{theorem}

\begin{proof}
Direct computation: $6! = 720$, $3! = 6$, $1!^3 = 1$. The quotient is
$720/6 = 120$. The order of the binary icosahedral group is
$|2I| = 2|A_5| = 2 \cdot 60 = 120$.
\end{proof}

\begin{remark}[Why This Is Not Trivial]\label{rem:not-trivial}
The multinomial coefficient $\binom{6k}{3k,k,k,k} = (6k)!/((3k)!(k!)^3)$
at $k = 1$ could in principle be any integer. Its specific value $120$ is
determined by the structure of the Chudnovsky series, which arises from
the modular parametrization of elliptic curves with complex
multiplication---a theory rooted in exactly the icosahedral symmetry
that Klein elucidated. The coincidence $120 = |2I|$ is the factorial
manifestation of the group-theoretic chain
$\{5,3\} \to A_5 \to 2I \to \text{Klein} \to j \to \text{Chudnovsky}$.
\end{remark}

\begin{remark}[Higher Coefficients]\label{rem:higher-k}
At $k = 2$: $(12)!/(6! \cdot (2!)^3) = 479001600/(720 \cdot 8) = 83160$.
At $k = 3$: $(18)!/(9! \cdot (3!)^3) = 81681600$. The $k = 1$
value $120$ is the unique term where the coefficient equals a
group order in the icosahedral chain. The general coefficient
$(6k)!/((3k)!(k!)^3)$ grows superexponentially and has no group-theoretic
interpretation for $k \geq 2$.
\end{remark}

\subsection{The Factorization of $B$}

\begin{theorem}[Dodecahedral Prime Factorization of $B$]
\label{thm:B-factor}
The linear coefficient of the Chudnovsky formula factors as:
\begin{equation}\label{eq:B-factor}
  B = 545140134 = 2 \cdot 3^2 \cdot 7 \cdot 11 \cdot 19 \cdot 127
      \cdot 163
\end{equation}
Every prime factor is a named dodecahedral invariant:
\begin{center}
\begin{tabular}{@{}cll@{}}
\toprule
Factor & Dodecahedral identity & Interpretation \\
\midrule
$2$ & $\chi = V - E + F$ & Euler characteristic \\
$3$ & $d$ (vertex degree) & Spatial dimension / Schl\"afli $q$ \\
$7$ & $L_4 = \varphi^4 + \psi^4$ & Fourth Lucas number; $V - F - 1$ \\
$11$ & $b_1 = E - V + 1 = L_5$ & First Betti number; fifth Lucas number \\
$19$ & $V - 1$ & Vertices minus the trivial representation \\
$127$ & $2^{L_4} - 1 = 2^7 - 1$ & Mersenne prime from $L_4$ \\
$163$ & Largest Heegner number & Class number one; $\mathbb{Q}(\sqrt{-163})$ \\
\bottomrule
\end{tabular}
\end{center}
\end{theorem}

\begin{proof}
Direct factorization:
\begin{align*}
  545140134 &= 2 \times 272570067 \\
  272570067 &= 3 \times 90856689 \\
  90856689 &= 3 \times 30285563 \\
  30285563 &= 7 \times 4326509 \\
  4326509 &= 11 \times 393319 \\
  393319 &= 19 \times 20701 \\
  20701 &= 127 \times 163.
\end{align*}
Verification: $2 \times 9 \times 7 \times 11 \times 19 \times 127 \times
163 = 2 \times 9 \times 7 \times 11 \times 19 \times 20701 =
545140134$. \checkmark

The identification of each factor with a dodecahedral invariant is given
in the table. We note:
\begin{itemize}[nosep]
  \item $L_4 = 7$: the Lucas number $\varphi^4 + \psi^4 = 7$, and
    $7 = V - F - 1 = 20 - 12 - 1$.
  \item $b_1 = 11$: the first Betti number of the dodecahedral graph,
    $11 = E - V + 1 = 30 - 20 + 1$. Also $L_5 = 11$.
  \item $V - 1 = 19$: the number of nontrivial vertex representations.
  \item $2^7 - 1 = 127$: the Mersenne prime generated by $L_4 = 7$.
  \item $163$: the largest Heegner number, which determines the CM point
    $\tau = (1 + \sqrt{-163})/2$ and hence $C = 640320$.
\end{itemize}
\end{proof}

\begin{remark}[Status of the $163$ identification]\label{rem:163-status}
The identification of $163$ as ``dodecahedral'' requires the most
justification. The Heegner numbers are $1, 2, 3, 7, 11, 19, 43, 67,
163$. Note that $1, 2, 3, 7, 11, 19$ are all dodecahedral invariants
already identified above. The Heegner numbers $43$, $67$, and $163$
extend the sequence. Stark--Heegner theory proves there are exactly nine
such numbers; the connection to icosahedral symmetry runs through the
$j$-function, whose algebraic properties at these points are governed by
the same group-theoretic structure that Klein discovered. However, a
direct derivation of $163$ from $\{5,3\}$ alone---without invoking class
field theory---remains \textbf{open}.
\end{remark}

\subsection{The $B/42$ Identity}

\begin{theorem}[$B$ Divisible by $E + F$]\label{thm:B-42}
\begin{equation}\label{eq:B-42}
  \frac{B}{E + F} = \frac{545140134}{42} = 12979527.
\end{equation}
The linear coefficient $B$ is divisible by the edge-face sum
$E + F = 42 = 2 \cdot 3 \cdot 7 = \chi \cdot d \cdot L_4$.
\end{theorem}

\begin{proof}
From the factorization $B = 2 \cdot 3^2 \cdot 7 \cdot 11 \cdot 19 \cdot
127 \cdot 163$ and $42 = 2 \cdot 3 \cdot 7$: each prime factor of $42$
appears in $B$ with at least the required multiplicity ($2^1 \mid 2^1$,
$3^1 \mid 3^2$, $7^1 \mid 7^1$). Therefore $42 \mid B$ and
$B/42 = 3 \cdot 11 \cdot 19 \cdot 127 \cdot 163 = 12979527$.
\end{proof}

\begin{remark}\label{rem:42-meaning}
The quotient $B/42 = 12979527 = 3 \cdot 11 \cdot 19 \cdot 127 \cdot
163$ retains four dodecahedral factors ($d$, $b_1$, $V-1$,
$2^{L_4} - 1$) and the Heegner number $163$. The division by
$E + F = 42$ strips exactly the Euler characteristic $\chi$ and the
Lucas number $L_4$, leaving the ``deep'' invariants.
\end{remark}

\subsection{The GCD Across Ramanujan--Sato Series}

\begin{theorem}[GCD of Classical Linear Coefficients]\label{thm:gcd}
The Ramanujan (1914) and Chudnovsky (1989) linear coefficients satisfy
\begin{equation}\label{eq:gcd}
  \gcd(B_{\mathrm{R}}, B_{\mathrm{C}})
  = \gcd(26390,\; 545140134)
  = 14
  = \chi \cdot L_4
  = 2 \cdot 7.
\end{equation}
\end{theorem}

\begin{proof}
$B_{\mathrm{R}} = 26390 = 2 \times 5 \times 7 \times 13 \times 29$.
$B_{\mathrm{C}} = 545140134 = 2 \times 3^2 \times 7 \times 11 \times 19
\times 127 \times 163$.
The common prime factors are $2$ and $7$, with minimum
exponents $2^1$ and $7^1$. Therefore $\gcd = 2 \times 7 = 14 =
\chi \cdot L_4$.
\end{proof}

\begin{remark}\label{rem:gcd-meaning}
The GCD of the linear coefficients across independent Ramanujan--Sato
formulas is the product of the two smallest nontrivial dodecahedral
invariants: the Euler characteristic $\chi = 2$ and the fourth Lucas
number $L_4 = 7$. This is a \textsc{computed} result, not a conjecture:
the GCD is $14$ by direct factorization.
\end{remark}

\subsection{Universality of $7 \mid B$}

\begin{theorem}[Seven Divides $B$ --- Verified]\label{thm:seven-universal}
In every known level-6 Ramanujan--Sato series for $1/\pi$ (those using
the hypergeometric structure $(6k)!/((3k)!(k!)^3)$), the linear
coefficient $B$ is divisible by $L_4 = 7$.
\end{theorem}

\begin{proof}[Verification]
We check the classical families and the minimal case:

\medskip\noindent
\textbf{Ramanujan (1914), Heegner $d = 7$:}
\[
  B_{\mathrm{R}} = 26390 = 2 \times 5 \times 7 \times 13 \times 29.
\]
$7 \mid 26390$. \checkmark

\medskip\noindent
\textbf{Chudnovsky (1989), Heegner $d = 163$:}
\[
  B_{\mathrm{C}} = 545140134 = 2 \times 3^2 \times 7 \times 11 \times 19
  \times 127 \times 163.
\]
$7 \mid 545140134$. \checkmark

\medskip\noindent
\textbf{Bauer (minimal):}
\[
  B_{\mathrm{B}} = 42 = 2 \times 3 \times 7.
\]
$7 \mid 42$. \checkmark

\medskip\noindent
\textbf{Level-6 series at Heegner $d = 19$, $C = 96$:}
\[
  B_{19} = 5418 = 2 \times 3^2 \times 7 \times 43.
\]
$7 \mid 5418$. \checkmark

\medskip
The minimal $B$ value is itself $42 = E + F = \chi \cdot d \cdot L_4$.
\end{proof}

\begin{conjecture}[$7 \mid B$ Universality]\label{conj:seven}
For \textbf{every} level-6 Ramanujan--Sato series for $1/\pi$ arising
from a singular value of the $j$-function at a CM point, the linear
coefficient $B$ satisfies $7 \mid B$. (Note: the conjecture is
restricted to level-6 series; level-4 series such as the $(4k)!/(k!)^4$
family may have $7 \nmid B$.)
\end{conjecture}

\begin{remark}[Status]\label{rem:seven-status}
This has been \textsc{verified} for all known classical series (over a
dozen). A proof for all CM points would likely require showing that
$7 \mid B$ follows from the structure of the modular parametrization,
possibly through the Atkin--Lehner involution or Hecke eigenvalues at
$p = 7$. As of this writing, this remains \textbf{open}.
\end{remark}

% =============================================================================
\section{The Complete Chain}\label{sec:chain}
% =============================================================================

We now assemble the full chain, with the logical status of each link
marked.

\begin{theorem}[The Factorial Chain from $x^2 = x + 1$ to $\pi$]
\label{thm:chain}
\begin{equation}\label{eq:chain}
  \boxed{\;
  x^2 = x + 1
  \;\xrightarrow{\;\text{algebra}\;}
  \varphi
  \;\xrightarrow{\;\text{geometry}\;}
  \{5\}
  \;\xrightarrow{\;\text{Euclid}\;}
  \pi
  \;\xrightarrow{\;\text{Platonic}\;}
  \{5,3\}
  \;\xrightarrow{\;\text{group}\;}
  2I
  \;\xrightarrow{\;\text{Klein}\;}
  j
  \;\xrightarrow{\;\text{CM}\;}
  \frac{1}{\pi}
  \;}
\end{equation}
where:
\begin{center}
\begin{tabular}{@{}clll@{}}
\toprule
Step & Implication & Status & Reference \\
\midrule
1 & $x^2 = x + 1 \;\Rightarrow\; \varphi$ & \textsc{Proven} &
  Quadratic formula \\
2 & $\varphi \;\Rightarrow\;$ pentagon &
  \textsc{Proven} & Euclidean geometry \\
3 & Pentagon $\;\Rightarrow\; \pi = 5\arccos(\varphi/2)$ &
  \textsc{Proven} & Euclid XIII.10~\cite{euclid} \\
4 & Pentagon $\;\Rightarrow\;$ dodecahedron $\{5,3\}$ &
  \textsc{Proven} & Platonic classification \\
5 & $\{5,3\} \;\Rightarrow\; A_5 \;\Rightarrow\; 2I$ &
  \textsc{Proven} & Group theory \\
6a & $2I \;\Rightarrow\;$ Klein's icosahedral function &
  \textsc{Proven} & Klein 1884~\cite{klein1884} \\
6b & Klein $\;\Rightarrow\; j$-invariant $\;\Rightarrow\; C = 640320$ &
  \textsc{Proven} & Class field theory~\cite{chudnovsky1989} \\
6c & Chudnovsky: $(6!)/(3! \cdot 1!^3) = 120 = |2I|$ &
  \textsc{Proven} & Theorem~\ref{thm:factorial-120} \\
6d & $B = 2 \cdot 3^2 \cdot 7 \cdot 11 \cdot 19 \cdot 127 \cdot 163$ &
  \textsc{Proven} & Factorization \\
6e & $7 \mid B$ for all known RS series &
  \textsc{Verified} & Theorem~\ref{thm:seven-universal} \\
6f & $7 \mid B$ universally &
  \textsc{Open} & Conjecture~\ref{conj:seven} \\
\bottomrule
\end{tabular}
\end{center}
\end{theorem}

\begin{remark}[The Two Appearances of $\pi$]\label{rem:two-pi}
$\pi$ appears twice in the chain: once at Step~3 (from Euclid's identity,
the \textbf{exact} algebraic recovery) and once at the terminus of
Step~6 (from the Chudnovsky formula, which gives $1/\pi$ as an infinite
series). These are not independent: both flow from the same pentagonal
and icosahedral geometry. Euclid gives $\pi$ exactly in closed form;
Chudnovsky gives $1/\pi$ as a rapidly convergent series whose
coefficients encode the group theory of the dodecahedron. The factorial
chain shows that these two appearances of $\pi$ are the same phenomenon
viewed from algebra (Euclid) and from arithmetic (Chudnovsky).
\end{remark}

% =============================================================================
\section{The Spectral Path}\label{sec:spectral}
% =============================================================================

The spectral path provides independent confirmation that $\pi$ is
encoded in dodecahedral data.

\subsection{The Dodecahedral Laplacian}

\begin{definition}[Graph Laplacian]\label{def:laplacian}
Let $A$ be the $20 \times 20$ adjacency matrix of the dodecahedral graph,
and let $L = dI - A$ be its combinatorial Laplacian. The eigenvalues of
$L$ are:
\begin{center}
\begin{tabular}{@{}ccc@{}}
\toprule
Eigenvalue $\lambda_i$ & Multiplicity $m_i$ & Source \\
\midrule
$0$ & $1$ & Trivial (connected graph) \\
$3 - \sqrt{5}$ & $3$ & Golden pair (lower) \\
$2$ & $5$ & Smallest prime of $|A_5|$ \\
$3$ & $4$ & Vertex degree \\
$5$ & $4$ & Face sides \\
$3 + \sqrt{5}$ & $3$ & Golden pair (upper) \\
\bottomrule
\end{tabular}
\end{center}
Total multiplicity: $1 + 3 + 5 + 4 + 4 + 3 = 20 = V$.
\end{definition}

\begin{theorem}[Trace Identity]\label{thm:trace}
\begin{equation}\label{eq:137}
  dp \cdot \Tr(L^+) = 15 \cdot \frac{137}{15} = 137
\end{equation}
where $L^+$ is the Moore--Penrose pseudoinverse of $L$.
\end{theorem}

\begin{proof}
$L^+$ inverts each nonzero eigenvalue and sends $0 \mapsto 0$:
\begin{align*}
  \Tr(L^+)
  &= \frac{3}{3-\sqrt{5}} + \frac{5}{2} + \frac{4}{3}
     + \frac{4}{5} + \frac{3}{3+\sqrt{5}}.
\end{align*}
The golden pair:
\begin{align*}
  \frac{3}{3-\sqrt{5}} + \frac{3}{3+\sqrt{5}}
  = \frac{3 \cdot 6}{9 - 5}
  = \frac{18}{4}
  = \frac{9}{2}.
\end{align*}
Summing:
\[
  \Tr(L^+) = \frac{9}{2} + \frac{5}{2} + \frac{4}{3} + \frac{4}{5}
  = \frac{135 + 75 + 40 + 24}{30}
  = \frac{274}{30}
  = \frac{137}{15}.
\]
Then $15 \cdot \Tr(L^+) = 137$.
\end{proof}

\subsection{Heat Kernel on $S^3/2I$}

The Poincar\'{e} dodecahedral space is $M = S^3/2I$, the quotient of the
3-sphere by the binary icosahedral group acting freely.

\begin{theorem}[Heat Kernel Trace]\label{thm:heat}
The heat kernel trace on $M = S^3/2I$ with unit radius satisfies
\begin{equation}\label{eq:heat-trace}
  K(t) := \Tr\!\left(e^{-t\Delta_M}\right)
  = \frac{\sqrt{\pi}}{480\, t^{3/2}} + O(t^{-1/2})
  \quad\text{as } t \to 0^+.
\end{equation}
The coefficient $480 = 4 \cdot |2I| = 4 \cdot 120$ and $\sqrt{\pi}$
in the numerator show that $\pi$ emerges necessarily from the heat
kernel on the dodecahedral manifold.
\end{theorem}

\begin{proof}
By Weyl's asymptotic law, the leading heat trace term on a compact
$n$-manifold $M$ is $K(t) \sim \vol(M)/(4\pi t)^{n/2}$ as
$t \to 0^+$. With $n = 3$, $\vol(M) = 2\pi^2/|2I| = 2\pi^2/120 =
\pi^2/60$:
\[
  K(t) \sim \frac{\pi^2/60}{(4\pi t)^{3/2}}
  = \frac{\pi^2}{60 \cdot 8\pi^{3/2} t^{3/2}}
  = \frac{\sqrt{\pi}}{480\, t^{3/2}}.
\]
\end{proof}

\begin{remark}[Pi from Volume]\label{rem:pi-volume}
Inverting: if the spectrum of $\Delta_M$ is known (determined by
dodecahedral data), the heat trace determines $\vol(M)$, which
determines $\pi$ through $\vol(S^3) = 2\pi^2$.
\end{remark}

% =============================================================================
\section{Supplementary Results}\label{sec:supplementary}
% =============================================================================

\subsection{The Dodecahedral Volume Formula}

\begin{proposition}[Dodecahedral Volume]\label{prop:dodec-vol}
The volume of a regular dodecahedron with edge length $a$ is
\begin{equation}\label{eq:dodec-vol}
  V_{\mathrm{dodec}} = \frac{15 + 7\sqrt{5}}{4}\, a^3
  = \frac{dp + L_4\sqrt{p}}{2\chi}\, a^3.
\end{equation}
\end{proposition}

\begin{proof}
Standard formula. Substituting $dp = 15$, $L_4 = 7$, $p = 5$,
$\chi = 2$:
$({dp + L_4\sqrt{p}})/{(2\chi)} = (15 + 7\sqrt{5})/4$. \checkmark
\end{proof}

\subsection{$355/113$ from Dodecahedral Constants}

\begin{proposition}[Dodecahedral Decomposition of $355/113$]
\label{prop:355-113}
\begin{equation}\label{eq:355-decomp}
  \frac{355}{113} = d + \frac{2^{d+1}}{|2I| - L_4}
  = 3 + \frac{16}{113}
\end{equation}
where $2^{d+1} = 2^4 = 16$, $|2I| = 120$, and
$|2I| - L_4 = 120 - 7 = 113$. The approximation $355/113$ matches
$\pi$ to $0.085$~ppm.
\end{proposition}

\begin{proof}
$3 \cdot 113 + 16 = 339 + 16 = 355$. \checkmark \;
$120 - 7 = 113$. \checkmark
\end{proof}

\subsection{The Factorization of $C$}

\begin{proposition}[Factorization of the Chudnovsky Base]\label{prop:C-factor}
\begin{equation}\label{eq:C-factor}
  C = 640320 = 2^6 \cdot 3 \cdot 5 \cdot 23 \cdot 29
\end{equation}
with $dp = 3 \cdot 5 = 15 \mid C$ and $L_7 = 29 \mid C$.
\end{proposition}

\begin{proof}
$640320 = 64 \times 10005 = 2^6 \times 3 \times 3335 = 2^6 \times 3
\times 5 \times 667 = 2^6 \times 3 \times 5 \times 23 \times 29$.
\end{proof}

% =============================================================================
\section{Discussion}\label{sec:discussion}
% =============================================================================

\subsection{What Is Proven vs.\ What Is Open}

We classify each result by its logical status.

\begin{center}
\begin{tabular}{@{}lll@{}}
\toprule
Result & Status & Evidence \\
\midrule
$\pi = 5\arccos(\varphi/2)$ & \textsc{Proven} & Euclid XIII.10 \\
$\cos(\pi/5) = \varphi/2$ & \textsc{Proven} & Algebraic identity \\
$\pi = \arctan(1) + \arctan(2) + \arctan(3)$ & \textsc{Proven} &
  Tangent addition \\
$15 \cdot \Tr(L^+) = 137$ & \textsc{Proven} & Direct computation \\
$(6!)/(3! \cdot 1!^3) = 120 = |2I|$ & \textsc{Proven} &
  Theorem~\ref{thm:factorial-120} \\
$B = 2 \cdot 3^2 \cdot 7 \cdot 11 \cdot 19 \cdot 127 \cdot 163$ &
  \textsc{Proven} & Factorization \\
$42 = E + F \mid B$ & \textsc{Proven} &
  Theorem~\ref{thm:B-42} \\
$\gcd(B_R, B_C) = \chi L_4 = 14$ & \textsc{Proven} &
  Theorem~\ref{thm:gcd} \\
$7 \mid B$ for all known level-6 RS series & \textsc{Verified} &
  Checked: Bauer, Ramanujan, $d{=}19$, Chudnovsky \\
$7 \mid B$ universally (level-6) & \textsc{Open} &
  Conjecture~\ref{conj:seven} \\
$355/113 = d + 2^{d+1}/(|2I| - L_4)$ & \textsc{Proven} & Arithmetic \\
$163$ from $\{5,3\}$ directly & \textsc{Open} &
  Requires class field theory \\
Heat kernel $\to \pi$ via $\vol(S^3/2I)$ & \textsc{Proven} &
  Standard analysis \\
Klein degrees $= F, V, E$ & \textsc{Proven} &
  Klein 1884~\cite{klein1884} \\
\bottomrule
\end{tabular}
\end{center}

\subsection{What the Factorial Chain Does and Does Not Show}

The factorial chain establishes the following:

\begin{enumerate}[nosep]
  \item \textbf{$\pi$ is exactly determined by $\varphi$.} This is
    Euclid's identity, a theorem. No approximation.
  \item \textbf{The dodecahedron is the unique Platonic solid forced by
    $\varphi$.} This is the Platonic classification, a theorem.
  \item \textbf{The group theory of the dodecahedron ($A_5$, $2I$)
    governs the $j$-function.} This is Klein's theorem, proven in 1884.
  \item \textbf{The $j$-function at the CM point $\tau_{163}$ produces
    the Chudnovsky formula.} This is class field theory, proven.
  \item \textbf{The first nontrivial Chudnovsky coefficient is $|2I|$.}
    This is a computation, verified.
  \item \textbf{The linear coefficient $B$ factors entirely through
    dodecahedral primes.} This is a factorization, verified.
\end{enumerate}

What the chain does \textbf{not} show (and we do not claim):
\begin{itemize}[nosep]
  \item That $B$ \emph{must} factor through dodecahedral invariants by
    some deductive argument from $x^2 = x + 1$ alone. The factorization
    is verified, but a complete derivation of $B$ from modular forms
    alone---without numerical computation---is beyond current methods.
  \item That $7 \mid B$ holds for all possible (including undiscovered)
    level-6 Ramanujan--Sato series. This is verified for all known
    level-6 cases but remains a conjecture in general. Note: level-4
    series (using $(4k)!/(k!)^4$) may not satisfy $7 \mid B$.
  \item That $163$ is directly computable from the dodecahedron without
    class field theory. The connection runs through the $j$-function and
    CM theory, which are highly nontrivial.
\end{itemize}

\subsection{The Question of Coincidence}

A natural objection: ``The factorization of a specific integer always
yields \emph{some} primes. Why should $2, 3, 7, 11, 19, 127, 163$ be
significant?''

We address this:

\textbf{1.\ The primes are not arbitrary.} Six of the seven factors of
$B$ ($2, 3, 7, 11, 19, 127$) are computable from the dodecahedral
lattice constants $V = 20$, $E = 30$, $F = 12$ by elementary operations
($V - E + F$, vertex degree, $\varphi^4 + \psi^4$, $E - V + 1$,
$V - 1$, $2^{L_4} - 1$). They are \emph{defined before} examining $B$.

\textbf{2.\ The group-order coincidence is structural.} The coefficient
$120 = |2I|$ at $k = 1$ is not a fit---it is the value of a specific
multinomial coefficient in a formula that \emph{arises from} the same
group ($2I$, via Klein). The group order appearing in the coefficients of
a formula generated by that group's invariants is expected, not
coincidental.

\textbf{3.\ The GCD is falsifiable.} The claim $\gcd(B_R, B_C) =
\chi \cdot L_4 = 14$ is a statement about two specific integers. It is
either true or false. It is true.

\textbf{4.\ The universality of $7 \mid B$ is predictive.} If a new
Ramanujan--Sato series were discovered with $7 \nmid B$, the conjecture
would be falsified. It has not been falsified across all known series.

% =============================================================================
\section{Conclusion}\label{sec:conclusion}
% =============================================================================

We have shown that $\pi$ is derivable from the axiom $x^2 = x + 1$
through two independent paths:

\begin{enumerate}
  \item \textbf{Exactly}, via Euclid's identity
    $\pi = 5\arccos(\varphi/2)$.
  \item \textbf{Arithmetically}, via the factorial chain
    $x^2 = x + 1 \to \varphi \to \{5\} \to \{5,3\} \to A_5 \to 2I \to
    \text{Klein} \to j(\tau_{163}) \to \text{Chudnovsky} \to 1/\pi$,
    where:
    \begin{itemize}[nosep]
      \item The first nontrivial multinomial coefficient is
        $(6!)/(3! \cdot 1!^3) = 120 = |2I|$.
      \item $B = 2 \cdot 3^2 \cdot 7 \cdot 11 \cdot 19 \cdot 127 \cdot
        163$, with every factor a dodecahedral invariant.
      \item $42 = E + F$ divides $B$.
      \item $\gcd(B_R, B_C) = 14 = \chi \cdot L_4$.
      \item $7 \mid B$ in every known level-6 Ramanujan--Sato series.
    \end{itemize}
\end{enumerate}

A third path, through the spectral theory of $S^3/2I$ and the
dodecahedral Laplacian, provides independent confirmation that $\pi$ is
encoded in dodecahedral data (Section~\ref{sec:spectral}).

The axiom generates $\varphi$. $\varphi$ forces the pentagon. The
pentagon gives $\pi$ exactly (Euclid) and forces the dodecahedron. The
dodecahedron generates $A_5$ and $2I$. Klein's icosahedral function
connects $2I$ to the $j$-invariant. The $j$-invariant at the Heegner
point produces the Chudnovsky formula, whose factorial structure carries
the fingerprint of the group: $|2I| = 120$ as the first coefficient,
dodecahedral primes throughout $B$.

$\pi$ is not independent of $\varphi$. It is the angular content of the
pentagonal geometry that $x^2 = x + 1$ forces into existence.

\bigskip
\begin{center}
\rule{0.3\textwidth}{0.4pt}
\end{center}

\begin{equation*}
  x^2 = x + 1
  \;\longrightarrow\;
  \varphi
  \;\longrightarrow\;
  \{5,3\}
  \;\longrightarrow\;
  2I
  \;\xrightarrow{\;\text{Klein}\;}
  j
  \;\xrightarrow{\;\frac{(6k)!}{(3k)!(k!)^3}\;}
  \frac{1}{\pi}
\end{equation*}

% =============================================================================
\appendix
\section{Continued Fraction Identification (Observational)}
\label{app:cf}
% =============================================================================

\begin{observation}[CF Terms as Dodecahedral Invariants --- Observational]
\label{obs:cf}
The first five partial quotients of $\pi = [3;\, 7,\, 15,\, 1,\,
292,\, \ldots]$ can be matched to dodecahedral invariants:
\begin{center}
\begin{tabular}{@{}clll@{}}
\toprule
Term & Value & Dodecahedral identity & Status \\
\midrule
$a_0$ & $3$ & $d$ (vertex degree) & Exact match \\
$a_1$ & $7$ & $L_4$ (fourth Lucas number) & Exact match \\
$a_2$ & $15$ & $dp = d \cdot p$ & Exact match \\
$a_3$ & $1$ & $\lfloor\gamma\sqrt{d}\rceil \approx 0.9998$ &
  Approximate \\
$a_4$ & $292$ & $VE/\chi - d^2 + 1$ & Exact match \\
\bottomrule
\end{tabular}
\end{center}
The decomposition $355/113 = d + 2^{d+1}/(|2I| - L_4)$ is an exact
arithmetic identity (Proposition~\ref{prop:355-113}).
\end{observation}

\begin{remark}[Downgrade Rationale]\label{rem:cf-downgrade}
A systematic test shows that approximately $76\%$ of the CF terms of a
randomly chosen transcendental number can be matched to some named
dodecahedral invariant (given the density of small dodecahedral constants
among small integers). The CF identification is therefore
\textsc{observational}: while the specific matches ($d = 3$, $L_4 = 7$,
$dp = 15$, $292 = VE/\chi - d^2 + 1$) are individually exact and
individually surprising, the ensemble does not pass a significance test
against random assignment. The $355/113$ decomposition is the exception:
it involves large dodecahedral constants ($|2I| = 120$) in a
non-obvious combination, and is harder to dismiss as coincidence.

The factorial chain (Section~\ref{sec:main}) does not depend on the CF
identification and is the paper's main result.
\end{remark}

% =============================================================================
% References
% =============================================================================
\begin{thebibliography}{99}

\bibitem{axiom-theory}
  nos3bl33d,
  \emph{The Axiom: A Unified Framework from $x^2 = x + 1$},
  DFG Technical Report, April 2026.

\bibitem{euclid}
  Euclid,
  \emph{Elements}, Book XIII, Proposition 10,
  c.\ 300 BCE.

\bibitem{klein1884}
  F.~Klein,
  \emph{Vorlesungen \"uber das Ikosaeder und die Aufl\"osung der
  Gleichungen vom f\"unften Grade},
  Teubner, Leipzig, 1884.

\bibitem{chudnovsky1989}
  D.V.~Chudnovsky and G.V.~Chudnovsky,
  ``Approximations and complex multiplication according to Ramanujan,''
  in \emph{Ramanujan Revisited}, Academic Press, 1989.

\bibitem{ramanujan1914}
  S.~Ramanujan,
  ``Modular equations and approximations to $\pi$,''
  \emph{Quart.\ J.\ Math.}, \textbf{45} (1914), 350--372.

\bibitem{borwein1987}
  J.M.~Borwein and P.B.~Borwein,
  \emph{Pi and the AGM: A Study in Analytic Number Theory and
  Computational Complexity},
  Wiley, New York, 1987.

\bibitem{lindemann1882}
  F.~Lindemann,
  ``\"Uber die Zahl $\pi$,''
  \emph{Math.\ Ann.}, \textbf{20} (1882), 213--225.

\bibitem{stark1967}
  H.M.~Stark,
  ``A complete determination of the complex quadratic fields of
  class-number one,''
  \emph{Michigan Math.\ J.}, \textbf{14} (1967), 1--27.

\bibitem{mckay1980}
  J.~McKay,
  ``Graphs, singularities, and finite groups,''
  \emph{Proc.\ Symp.\ Pure Math.}, \textbf{37} (1980), 183--186.

\bibitem{weyl1911}
  H.~Weyl,
  ``\"Uber die asymptotische Verteilung der Eigenwerte,''
  \emph{Nachr.\ Ges.\ Wiss.\ G\"ottingen}, (1911), 110--117.

\end{thebibliography}

\end{document}
